% GENERATED FILE. Edit canonical lessons, then run npm run generate:books.
% canonical-source-hash: fnv1a64:4b928b6503c85550
% canonical-lessons: AR-C168-dhira, AR-C168-jism, AR-C168-kahil, AR-C168-halq, AR-C168-suda

\chapter{Arm, Body, Ankle, Throat, Headache}
\label{ch:ar-arm-body-ankle-throat-headache}
\hlchaptermodality{168}

\begin{chapteropening}
\textbf{By the end of this chapter:} \emph{I can say an arm, the body, an ankle, the throat and a headache.}

Complete the last lesson of chapter 168: I can say an arm, the body, an ankle, the throat and a headache.
\end{chapteropening}

\section[\texorpdfstring{\ar{ذراع} (dhirāʿ)}{ذراع (dhirāʿ)}]{\ar{ذراع} (dhirāʿ) — an arm}

\label{lesson:AR-C168-dhira}



\begin{quote}
Before the new one: say the Arabic for to escape, to run away, then the Arabic for to disappear.
\end{quote}

\subsection*{You'll want to know: \ar{ذراع}}
\textbf{\ar{ذراع}} — \emph{dhirāʿ} — \textquotedblleft{}an arm\textquotedblright{}.

\begin{cousinweb}
\textbf{\ar{ذراع}} (\emph{dhirāʿ}) is built on the native Arabic root \emph{dh-r-ʿ}, \textquotedblleft{}the forearm; a cubit\textquotedblright{} (Lane). Hebrew \emph{zeroaʿ}, \textquotedblleft{}arm\textquotedblright{}, is its cousin.
\end{cousinweb}

\begin{grammarlens}[title={What you've built}]
One more word for this chapter.
\end{grammarlens}

\subsection*{Your turn}
\begin{itemize}
\raggedright
  \item \emph{Say it:} \emph{dhirāʿ}
  \item \emph{Say it:} \emph{dhirāʿ}, once more
  \item \emph{Say it:} say \emph{dhirāʿ}
  \item \emph{Recall:} say the Arabic for to escape, to run away, then the Arabic for to disappear, then say \emph{dhirāʿ} again
\end{itemize}

\begin{tcolorbox}[breakable,colback=teal!4,colframe=teal!35!black,title={Before you move on}]
What does \ar{ذراع} mean? (\textquotedblleft{}An arm\textquotedblright{}.) Say it once more.
\end{tcolorbox}
\section[\texorpdfstring{\ar{جسم} (jism)}{جسم (jism)}]{\ar{جسم} (jism) — the body}

\label{lesson:AR-C168-jism}



\begin{quote}
Before the new one: say the Arabic for to disappear, then the Arabic for an arm.
\end{quote}

\subsection*{You'll want to know: \ar{جسم}}
\textbf{\ar{جسم}} — \emph{jism} — \textquotedblleft{}the body\textquotedblright{}.

\begin{cousinweb}
\textbf{\ar{جسم}} (\emph{jism}) is built on the native Arabic root \emph{j-s-m}, \textquotedblleft{}to be large of body\textquotedblright{} (Lane). \emph{Jism} is the body.
\end{cousinweb}

\begin{grammarlens}[title={What you've built}]
One more word for this chapter.
\end{grammarlens}

\subsection*{Your turn}
\begin{itemize}
\raggedright
  \item \emph{Say it:} \emph{jism}
  \item \emph{Say it:} \emph{jism}, once more
  \item \emph{Say it:} say \emph{jism}
  \item \emph{Recall:} say the Arabic for to disappear, then the Arabic for an arm, then say \emph{jism} again
\end{itemize}

\begin{tcolorbox}[breakable,colback=teal!4,colframe=teal!35!black,title={Before you move on}]
What does \ar{جسم} mean? (\textquotedblleft{}The body\textquotedblright{}.) Say it once more.
\end{tcolorbox}
\section[\texorpdfstring{\ar{كاحل} (kāḥil)}{كاحل (kāḥil)}]{\ar{كاحل} (kāḥil) — an ankle}

\label{lesson:AR-C168-kahil}



\begin{quote}
Before the new one: say the Arabic for an arm, then the Arabic for the body.
\end{quote}

\subsection*{You'll want to know: \ar{كاحل}}
\textbf{\ar{كاحل}} — \emph{kāḥil} — \textquotedblleft{}an ankle\textquotedblright{}.

\begin{cousinweb}
\textbf{\ar{كاحل}} (\emph{kāḥil}) has no entry in Lane. It shares its letters with \emph{k-ḥ-l}, \textquotedblleft{}to blacken the eyelids with kohl\textquotedblright{} (Lane), the root of \emph{kuḥl}, \textquotedblleft{}kohl\textquotedblright{}, and of English \emph{alcohol}, from \emph{al-kuḥl}; how the ankle came by the word is not settled.
\end{cousinweb}

\begin{grammarlens}[title={What you've built}]
One more word for this chapter.
\end{grammarlens}

\subsection*{Your turn}
\begin{itemize}
\raggedright
  \item \emph{Say it:} \emph{kāḥil}
  \item \emph{Say it:} \emph{kāḥil}, once more
  \item \emph{Say it:} say \emph{kāḥil}
  \item \emph{Recall:} say the Arabic for an arm, then the Arabic for the body, then say \emph{kāḥil} again
\end{itemize}

\begin{tcolorbox}[breakable,colback=teal!4,colframe=teal!35!black,title={Before you move on}]
What does \ar{كاحل} mean? (\textquotedblleft{}An ankle\textquotedblright{}.) Say it once more.
\end{tcolorbox}
\section[\texorpdfstring{\ar{حلق} (ḥalq)}{حلق (ḥalq)}]{\ar{حلق} (ḥalq) — the throat}

\label{lesson:AR-C168-halq}



\begin{quote}
Before the new one: say the Arabic for the body, then the Arabic for an ankle.
\end{quote}

\subsection*{You'll want to know: \ar{حلق}}
\textbf{\ar{حلق}} — \emph{ḥalq} — \textquotedblleft{}the throat\textquotedblright{}.

\begin{cousinweb}
\textbf{\ar{حلق}} (\emph{ḥalq}) is built on the native Arabic root \emph{ḥ-l-q}, \textquotedblleft{}the throat; to shave\textquotedblright{} (Lane). \emph{Ḥalq} is the throat.
\end{cousinweb}

\begin{grammarlens}[title={What you've built}]
One more word for this chapter.
\end{grammarlens}

\subsection*{Your turn}
\begin{itemize}
\raggedright
  \item \emph{Say it:} \emph{ḥalq}
  \item \emph{Say it:} \emph{ḥalq}, once more
  \item \emph{Say it:} say \emph{ḥalq}
  \item \emph{Recall:} say the Arabic for the body, then the Arabic for an ankle, then say \emph{ḥalq} again
\end{itemize}

\begin{tcolorbox}[breakable,colback=teal!4,colframe=teal!35!black,title={Before you move on}]
What does \ar{حلق} mean? (\textquotedblleft{}The throat\textquotedblright{}.) Say it once more.
\end{tcolorbox}
\section[\texorpdfstring{\ar{صداع} (ṣudāʿ)}{صداع (ṣudāʿ)}]{\ar{صداع} (ṣudāʿ) — a headache}

\label{lesson:AR-C168-suda}



\begin{quote}
Before the new one: say the Arabic for an ankle, then the Arabic for the throat.
\end{quote}

\subsection*{You'll want to know: \ar{صداع}}
\textbf{\ar{صداع}} — \emph{ṣudāʿ} — \textquotedblleft{}a headache\textquotedblright{}.

\begin{cousinweb}
\textbf{\ar{صداع}} (\emph{ṣudāʿ}) is built on the native Arabic root \emph{ṣ-d-ʿ}, \textquotedblleft{}to split\textquotedblright{} (Lane). \emph{Ṣudāʿ} has the \emph{fuʿāl} shape of illness words, like \emph{suʿāl}, \textquotedblleft{}a cough\textquotedblright{}.
\end{cousinweb}

\begin{grammarlens}[title={What you've built}]
One more word for this chapter.
\end{grammarlens}

\subsection*{Your turn}
\begin{itemize}
\raggedright
  \item \emph{Say it:} \emph{ṣudāʿ}
  \item \emph{Say it:} \emph{ṣudāʿ}, once more
  \item \emph{Say it:} say \emph{ṣudāʿ}
  \item \emph{Recall:} say the Arabic for an ankle, then the Arabic for the throat, then say \emph{ṣudāʿ} again
\end{itemize}

\begin{tcolorbox}[breakable,colback=teal!4,colframe=teal!35!black,title={Before you move on}]
What does \ar{صداع} mean? (\textquotedblleft{}A headache\textquotedblright{}.) Say it once more.
\end{tcolorbox}
